\subsubsection{Crossflow Drag}
Changing the crossflow coefficient affects aerodynamic braking, pitching moment, and rotational damping together. With the axial coefficient fixed at 0.8, we compare crossflow coefficients of 0.8, 1.2, and 1.6 using the same local optimization method and touchdown requirements.

\begin{table}[htbp]
\centering
\caption{Crossflow comparison. Inventories are assigned only to accepted candidates. The unresolved row is a search limitation, not proof that landing is physically impossible.}
\label{tab:p3-drag}
\begin{tabular}{@{}llrr@{}}
\toprule
$C_n$ & Outcome & Inventory (t) & Retained (t) \\
\midrule
0.8 & No accepted candidate & -- & -- \\
1.2 & Nominal reference & 26.228 & 5.078 \\
1.6 & Accepted local candidate & 27.130 & 5.072 \\
\bottomrule
\end{tabular}
\end{table}

At $C_n=1.6$, inventory increases by 0.901 t (3.44\%), while peak specific force decreases from approximately $7.703\,g_0$ to $7.187\,g_0$. This local result trades greater inventory for a lower peak load; larger crossflow drag does not automatically imply lower fuel demand. At $C_n=0.8$, the local solve reaches its 140 s limit without an accepted trajectory. We assign no fuel requirement to its failed candidate and do not infer a feasible propellant range across the full coefficient interval.
